\section{Related Work}
\label{sec:related}

\para{Copy mechanisms in sequence models.}
Pointer Networks~\cite{vinyals2015pointer} introduced the idea of using attention scores as a pointer to select elements from the input, enabling models to produce outputs that are not in a fixed vocabulary. \citet{gu2016copynet} proposed CopyNet, which uses a shared softmax over generate and copy modes, allowing the two to compete through a single normalization term. \citet{see2017get} introduced the pointer-generator network with an explicit generation probability $\pgen$ computed from decoder state, context vector, and input embedding. This $\pgen$ gate is the most widely used copy decision mechanism in practice. \citet{gulcehre2016pointing} proposed an MLP-based switching network to decide between pointing and generating, while \citet{merity2016pointer} used a sentinel mechanism where a learnable vector competes with pointer targets to determine the copy-generate mixture. All of these approaches embed the copy decision within the full generation model, sharing representations and training jointly. We isolate the copy decision into a separately trained module to test whether this specialization improves decision quality.

\para{Neural network calibration.}
\citet{guo2017calibration} demonstrated that modern deep networks are poorly calibrated---they tend to be overconfident---and that simple temperature scaling (a single learned parameter) can substantially improve calibration. Their work introduced expected calibration error (ECE) and reliability diagrams as standard evaluation tools. Subsequent work has extended calibration analysis to various domains, including natural language processing. However, no prior work has applied calibration analysis to the copy-or-generate decision in pointer-generator models. We bridge this gap by measuring ECE, maximum calibration error (MCE), and Brier scores for copy probabilities.

\para{Task specialization in neural networks.}
The question of whether specialized models outperform general-purpose models on specific subtasks has a long history. Mixture-of-experts architectures~\cite{vaswani2017attention} route inputs to specialized subnetworks, and multi-task learning research has shown that sharing representations across tasks can both help and hurt performance depending on task relatedness. Our work asks a specific version of this question: does a model trained only on the copy decision outperform one that learns copy decisions as a byproduct of the full generation task? Unlike prior work on task specialization, we focus specifically on \emph{calibration}---not just accuracy---of the specialized model's predictions.